\scanpage{110}
\uline{Cor.} Conditions équivalentes

(i) $D$ est anisotrope i.e. $\det_A\mid D$ est non nulle sur chaque fibre

(ii) $\delta(B)(=\delta(D))$ est un élémt inv. de $\deleted{\unsure{B^{\otimes-2}}}\inserted{(\det B)^{\otimes-2}}$

\deleted{i.e. $\det_{\unsure{A}}\mid D$ est partout non nulle}

i.e. $\det_B=\det_A\mid B$ est non dégénérée

(iii) $B$ n'est une sous-alg. sur aucune fibre

(iv) $B$ engendre $A$ comme \uline{algèbre} (pas comme module bien sûr)

\margin{\unsure{si} $2$ inv. sur $S$}
\[
\left\{\begin{aligned}
&(\mathrm{\unsure{v|iv}})\quad D\not\subset B\quad\text{sur tte fibre}\\
&(\mathrm{vi})\quad\varphi\mid B\text{ est non dégénérée}
\end{aligned}\right.
\]

Enfin, si $u,v$ sont deux éléments de $B$ tels que $1,u,v$ font une base, de $B$, ces conditions équivalent aussi à

(v) $1,u,v,uv$ forment une base de $A$

ou encore

(v$'$) $uv\bmod B$ est une \deleted{\unsure{section inv.}} \inserted{élémt inv.} de $A/B$

\begin{center}\rule{1.5cm}{0.35pt}\end{center}

\deleted{Supposons}

Nous regardons maintenant deux sous-fibrés \deleted{\unsure{vect.}} \deleted{\unsure{A}} $H,H'$ de rang $2$, contnt $1$ (ce sont néc. des sous-algèbres, par H.C.) et on suppose
\[
H\cap H'=k\qquad\text{sur tte fibre}
\]
donc
\[
H\oplus H'\overset{\text{déf.}}{=}B\qquad\text{sous-fibré de rang }3
\]
% archive sheet number: 109
